
\chapter{Literature Review}
\label{ch:literature}

\section{Deep learning for plant disease classification}

The dominant approach treats disease identification as image classification over a
fixed label set. Mohanty et al.\ demonstrated that a fine-tuned deep network reaches
high accuracy on PlantVillage and, in the same paper, that performance falls sharply
on images taken under different conditions \cite{mohanty2016}. That second result is
the one this project builds on, and it is frequently dropped when the first is
quoted.

Ferentinos \cite{ferentinos2018} extended coverage to a larger label set using
transfer learning from ImageNet weights, establishing the recipe that most later
work follows. Too et al.\ \cite{too2019} compared several modern architectures on
PlantVillage and found small differences between them, all above 99\,\%. The
consistency across architectures is itself informative: when four different model
families converge on the same number, the number is describing the dataset more than
the models.

\section{Datasets and the conditions they encode}

PlantVillage \cite{hughes2015} is the standard benchmark. Its images were collected
to isolate the leaf: the leaf is centred, the background is uniform, and lighting is
controlled. These choices make the classification task clean and make the resulting
model's assumptions narrow.

PlantDoc \cite{singh2020} was built to represent field conditions, with cluttered
backgrounds, variable lighting and multiple leaves. It has fewer images and fewer
classes, which limits how precisely it can be evaluated, and it is the closest
available proxy for the deployment setting.

The two datasets do not share a label space. PlantDoc contains classes PlantVillage
does not, and some PlantDoc images are ambiguous between categories. Any cross
dataset evaluation therefore requires an explicit mapping, and the mapping decision
changes the reported numbers. Chapter~\ref{ch:datasets} documents ours.

\section{Domain shift and shortcut learning}

The observation that models exploit incidental features correlated with labels in
the training set is well established outside agriculture, notably in Geirhos et
al.\ \cite{geirhos2020}, who show that standard training encourages reliance on
texture and context over shape. In the plant setting, the incidental feature
candidate is the background: soil, hands, pots and laboratory surfaces correlate
with disease category in curated data because collection was organised by category.

The mechanism story that follows is usually stated as a hypothesis and rarely
tested. If background is the shortcut, then randomising background during training
should reduce the field gap. That is a testable prediction, and to our knowledge it
has not been tested with a pre-registered direction and a seed-level paired design
in this setting.

\section{Explainability}

Grad-CAM \cite{selvaraju2017} is the standard tool for localising the evidence
behind a convolutional prediction. It is widely used in plant disease work to argue
that the model is looking at the lesion.

A heat map is an observation about gradient magnitude, not a measurement of causal
dependence. A model can attend to a lesion and still classify correctly from
background context, and a diffuse map can indicate uncertainty as easily as
attention. This project therefore treats Grad-CAM as a qualitative illustration and
uses an intervention design to answer the causal question.

\section{Calibration and selective prediction}

A classifier that reports a probability is only useful if that probability means
something. Guo et al.\ \cite{guo2017} showed that modern networks are often
overconfident, and that a temperature scaling parameter estimated on held-out data
corrects much of the error. Selective prediction, in which the model abstains below
a confidence threshold, converts a probability into an operating point: a coverage
level and the accuracy retained at that coverage.

For a screening tool the relevant question is not average accuracy but retained
accuracy at usable coverage. This project measures that directly rather than
assuming a threshold works.

\section{Summary of the gap in prior work}

\begin{table}[H]
\centering
\caption{What prior work establishes and what it leaves open}
\label{tab:litsummary}
\tabfont
\begin{tabularx}{\textwidth}{@{}p{3.1cm}Y Y@{}}
\toprule
\textbf{Area} & \textbf{Established} & \textbf{Left open} \\
\midrule
Curated accuracy & Above 99\,\% on PlantVillage across many architectures & Whether the number transfers \\
Field performance & Falls sharply on field images & How far, measured with intervals \\
Mechanism & Background suspected as shortcut & Whether background is relied on, measured by intervention \\
Remediation & Augmentation proposed & Whether it works, tested with seeds and a pre-registered direction \\
Usefulness & Calibration discussed & Retained accuracy at usable coverage \\
\bottomrule
\end{tabularx}
\end{table}

The contribution of this project sits in the four rows that say ``left open''. Each
becomes a measured quantity rather than a claim.
